\documentclass{article}
\usepackage[a4paper, total={6.5in, 10in}]{geometry}
\setlength{\parindent}{0pt}
\usepackage{tcolorbox}
\tcbset{colback=white!94!black, colframe=black!60!white, coltitle = white, boxrule=0.8pt, arc=2mm}
\newtcolorbox{boxs}[2][]{title= #2,#1}
\usepackage{datetime2}
\usepackage[british]{babel}
\usepackage{amsfonts}
\usepackage{amsmath}

\title{\textbf{Linear Algebra HW 5}}
\author{Robert O'Connell}
\date{\today}
\begin{document}
\maketitle
\vspace{5mm}


\subsection*{Question 1}

\[
D: \mathbb{P}_3[x] \rightarrow \mathbb{P}_2[x]
\]
\[
D(p(x)) = p  '(x) \quad \forall \ p(x) \in \mathbb{P}_3[x]
\]

Let \(p(x) \in \mathbb{P}_3[x]\)

Then \(\ \ p(x) = ax^3 + bx^2 + cx +d, \quad p'(x) = 3ax^2 + 2bx + c\)
\\

The matrix we see is defined from the coordinate mappings of \(\mathbb{P}_3[x]\) to the 
coordinate mapping of \(\mathbb{P}_2[x]\). We need to find the bases of these spaces such 
that the matrix representation from one coordinate mapping to the other has the listed form.
\\

We see that in the differentiation that all coeffiecents are taken down a degree and multiplied
by their previous degree.
\\
If we choose a standard basis for \(\mathbb{P}_3[x]\), \(\{x^3, x^2, x, 1\}\), it's coordinates are
\(\begin{bmatrix}
    a \\ b \\ c \\ d
\end{bmatrix}\)
\\

We then need the coordinate representation of \(p'(x) = 3ax^2 + 2bx + c\) with respect to the basis 
of \(\mathbb{P}_2[x]\) to be \(\begin{bmatrix}
    a \\ b \\ c
\end{bmatrix}\), since the given matrix does not scale any component, it only shifts each component 
down a degree

The basis \(\{3x^2, 2x , 1\}\) will achieve this, since this represents \(p'(x)\) as 
\(\begin{bmatrix}
    a \\ b \\ c
\end{bmatrix}\)




\subsection*{Question 2}

If we consider the coordinate mapping of the polynomials in \(\mathbb{P}_2[x]\) with respect to 
the basis \(\{x^2,x,1\}\), this now allows us to consider this problem as a mapping from 
\(\mathbb{R}^2 \rightarrow \mathbb{R}^3\).
\\


Every \(\begin{bmatrix}
     a \\ b
\end{bmatrix}\) can be written as a linear combination of \(\begin{bmatrix}
-3 \\ 1
\end{bmatrix}, \begin{bmatrix}
8 \\ -3
\end{bmatrix}\) (two linearly independent vectors in \(\mathbb{R}^2\)).

\begin{align*}
    \begin{bmatrix}
        a \\ b
    \end{bmatrix} &= 
    c_1 \begin{bmatrix}
        -3 \\ 1
    \end{bmatrix} + c_2 \begin{bmatrix}
        8 \\ -3
    \end{bmatrix} \\
    &= \begin{bmatrix}
        -3 & 8 \\
        1 & -3
    \end{bmatrix} 
    \begin{bmatrix}
        c_1 \\ c_2
    \end{bmatrix} \\
    \begin{bmatrix}
        -3 & -8 \\
        -1 & -3
    \end{bmatrix}
    \begin{bmatrix}
        a \\ b
    \end{bmatrix} &= \begin{bmatrix}
        c_1 \\ c_2
    \end{bmatrix} \\
    \begin{bmatrix}
        c_1 \\ c_2
    \end{bmatrix} &= 
    \begin{bmatrix}
        -3a - 8b \\
        -a -3b
    \end{bmatrix}
\end{align*}

\begin{align*}
    \Rightarrow T\begin{bmatrix}
    a \\ b
    \end{bmatrix} &= c_1 T \begin{bmatrix}
        -3 \\ 1
    \end{bmatrix} + c_2 T \begin{bmatrix}
        8 \\ -3
    \end{bmatrix} \\
    &= (-3a -8b)\begin{bmatrix}
        5 \\ -2 \\ 1
    \end{bmatrix} + (-a -3b)\begin{bmatrix}
        -3 \\ 8 \\ -2
    \end{bmatrix} \\
    &= \begin{bmatrix}
        -12a -31b \\
        -2a -8b \\
        -5a -10b
    \end{bmatrix}
\end{align*}

Putting this result back into polynomial form, which is easy since our basis is a simple one, we get
\[
T \begin{bmatrix}
    a \\ b
\end{bmatrix} = (-12a-31b)x^2 + (-2a -8b)x + (-5a -10b)
\]

\subsection*{Question 3}

Recall that the columns of the matrix \([T]_{\mathcal{B}, \mathcal{C}}\) has as its columns the images
of the basis vectors \(\mathcal{B}\) written with respect to \(\mathcal{C}\)

\[
[T]_{\mathcal{B}, \mathcal{C}} = \begin{bmatrix}
    [T(1)]_\mathcal{C} & [T(1 + x)]_\mathcal{C} & [T(1+x+x^2)]_\mathcal{C} & [T(1+x+x^2+x^3)]_\mathcal{C}
\end{bmatrix}
\]

We then calculate this images using the definition of the transformation

\[
T(1) = 0
\]
\[
T(1+x) = x
\]
\begin{align*}
    T(1+x+x^2) &= x + 2x^2 -10 \\
    &= 2x^2 + x -10
\end{align*}
\begin{align*}
    T(1+x+x^2+x^3) &= x + 2x^2 - 10 + 3x^3-30 x \\
    &= 3x^3 + 2x^2 -29x -10
\end{align*}

Representing these images with respect to the basis \(\mathcal{C}\) is simple

Thus,
\[
[T]_{\mathcal{B}, \mathcal{C}} = \begin{bmatrix}
    0 & 0 & -10 & -10 \\
    0 & 1 & 1 & -29 \\
    0 & 0 & 2 & 2 \\
    0 & 0 & 0 & 3
\end{bmatrix}
\]


\subsection*{Question 4}
\subsubsection*{(a)}
In asking what is the coordinate vector of \(\vec{w}\) with respect to \(\mathcal{F}\), we are 
asking for what scalar multiplied by \(\mathcal{F}\) give \(\vec{w}\)

Thus, we are looking for the constants \(c_1,c_2\) such that
\begin{align*}
    \begin{bmatrix}
        29 \\ -66
    \end{bmatrix} &= c_1\begin{bmatrix}
        -1 \\ 2
    \end{bmatrix} + c_2 \begin{bmatrix}
        -4 \\ 9
    \end{bmatrix} \\
    \begin{bmatrix}
        29 \\ -66
    \end{bmatrix} &= \begin{bmatrix}
        -1 & -4 \\
        2 & 9 \\
    \end{bmatrix} \begin{bmatrix}
        c_1 \\ c_2
    \end{bmatrix} \\
    \begin{bmatrix}
        -9 & -4 \\
        2 & 1
    \end{bmatrix} \begin{bmatrix}
        29 \\ -66
    \end{bmatrix} &= \begin{bmatrix}
        c_1 \\ c_2
    \end{bmatrix} \\
    \begin{bmatrix}
        c_1 \\ c_2
    \end{bmatrix} &= \begin{bmatrix}
        3 \\-8
    \end{bmatrix}
\end{align*}
Thus, \([w]_\mathcal{F} = \begin{bmatrix}
    3 \\ -8
\end{bmatrix}\)


\subsubsection*{(b)}

We know that the columns of \(P_{\mathcal{F}\rightarrow \mathcal{E}}\) are the vectors of 
\(\mathcal{F}\) writen with respect to \(\mathcal{E}\)

Thus,
\[
P_{\mathcal{F}\rightarrow \mathcal{C}} = \begin{bmatrix}
    -1 & -4 \\
    2 & 9
\end{bmatrix}
\]

To verify we compute \(P_{\mathcal{F}\rightarrow 
\mathcal{E}}[\vec{w}]_\mathcal{F} = [\vec{w}]_\mathcal{E}\), which should give 
\(\begin{bmatrix}
    29 \\ -66
\end{bmatrix}\)

\[
\begin{bmatrix}
    -1 & -4 \\
    2 & 9 
\end{bmatrix} \begin{bmatrix}
    3 \\ -8
\end{bmatrix} = \begin{bmatrix}
    29 \\ -66
\end{bmatrix}
\]
which it does.

\subsubsection*{(c)}

For this we will use the change of basis matrix \(P_{\mathcal{E}\rightarrow \mathcal{F}}\),
which we computed earlier in (a), note this is the inverse of the matrix 
\(P_{\mathcal{E}\rightarrow \mathcal{F}}\).
\[
P_{\mathcal{E}\rightarrow \mathcal{F}} = \begin{bmatrix}
        -9 & -4 \\
        2 & 1
    \end{bmatrix}
\]

We can now multiply the vectors \(\vec{e}_1, \vec{e}_2\) by this matrix to get the result.

\[
[\vec{e}_1]_\mathcal{F} = \begin{bmatrix}
        -9 & -4 \\
        2 & 1
    \end{bmatrix} \begin{bmatrix}
        1 \\ 0
    \end{bmatrix} = \begin{bmatrix}
        -9 \\ 2
    \end{bmatrix}
\]
\[
[\vec{e}_2]_\mathcal{F} = \begin{bmatrix}
        -9 & -4 \\
        2 & 1
    \end{bmatrix} \begin{bmatrix}
        0 \\ 1
    \end{bmatrix} = \begin{bmatrix}
        -4 \\ 1
    \end{bmatrix}
\]

\subsubsection*{(d)}
We have already found this to be
\[
P_{\mathcal{E}\rightarrow \mathcal{F}} = \begin{bmatrix}
        -9 & -4 \\ 
        2 & 1
    \end{bmatrix}
\]

To verify we compute 
\[
\begin{bmatrix}
        -9 & -4 \\
        2 & 1
    \end{bmatrix} \begin{bmatrix}
        -1 & -4 \\
        2 & 9
    \end{bmatrix} = \begin{bmatrix}
        1 & 0 \\
        0 & 1
    \end{bmatrix}
\]
which holds.


\subsubsection*{(e)}
We have that 
\(
[T]_{\mathcal{F}, \mathcal{E}} = \begin{bmatrix}
    6 & -2 \\
    -3 & -9
\end{bmatrix}
\)
since we are given the images of \(\mathcal{F}\) 
under \(T\) with respect to \(\mathcal{E}\), which form the columns of this matrix

Since our linear transformation takes in vectors with basis \(\mathcal{E}\) to find the matrix
\([T]_{\mathcal{E}, \mathcal{E}}\), we need to find where the basis vectors of \(\mathcal{E}\)
go under the transformation \(T\), i.e. \(T(\vec{e}_1), T{(\vec{e}_2)}\), and these will
make up the columns of the matrix \([T]_{\mathcal{E},\mathcal{E}}\)
\\

At the moment our transformation only takes in vectors with basis \(\mathcal{F}\), so what we can 
do is use what was calculated in (c), the basis vectors of \(\mathcal{E}\) with respect to 
\(\mathcal{F}\), and put them through \(T\), giving us what we wanted.

\[
T(\vec{e}_1) = \begin{bmatrix}
    6 & -2 \\
    -3 & -9
\end{bmatrix} \begin{bmatrix}
    -9 \\2
\end{bmatrix} = \begin{bmatrix}
    -58 \\ 9
\end{bmatrix}
\]
\[
T(\vec{e}_2) = \begin{bmatrix}
    6 & -2 \\
    -3 & -9
\end{bmatrix} \begin{bmatrix}
    -4 \\ 1
\end{bmatrix} = \begin{bmatrix}
    -26 \\ 3
\end{bmatrix} 
\]

Thus,
\[
[T]_{\mathcal{E},\mathcal{E}} = \begin{bmatrix}
        -58 & -26 \\
        9 & 3
    \end{bmatrix}
\]



\subsubsection*{(f)}

We have that 
\(
[T]_{\mathcal{F}, \mathcal{E}} = \begin{bmatrix}
    6 & -2 \\
    -3 & -9
\end{bmatrix}
\)
\\

We are checking whether \([T]_{\mathcal{E},\mathcal{E}} = 
[T]_{\mathcal{F},\mathcal{E}}P_{\mathcal{E} \rightarrow \mathcal{F}}\)

\begin{align*}
    [T]_{\mathcal{E},\mathcal{E}} &= \begin{bmatrix}
        6 & -2 \\
        -3 & -9
    \end{bmatrix} \begin{bmatrix}
            -9 & -4 \\
            2 & 1
        \end{bmatrix} \\
    &= \begin{bmatrix}
        -58 & -26 \\
        9 & 3
    \end{bmatrix}
\end{align*}
which holds.


\subsection*{Question 5}
\subsubsection*{(a)}
To show that \(T\) is a linear transformation we must show that it preserves addition
and scalar multiplication.
\\

Let \(p_1(x), p_2(x) \in \mathbb{P}_2[x]\) and let \(\lambda \in \mathbb{R}\)
\begin{itemize}
    \item \(T(p_1(x) + p_2(x)) = \begin{bmatrix}
        p_1(-1) + p_2(-1) \\
        p_1(0) + p_2(0) \\
        p_1(1) + p_2(1) 
    \end{bmatrix} = \begin{bmatrix}
        p_1(-1)\\
        p_1(0) \\
        p_1(1)
    \end{bmatrix} + \begin{bmatrix}
        p_2(-1)\\
        p_2(0) \\
        p_2(1)
    \end{bmatrix} = T(p_1(x)) + T(p_2(x))\)

    \item \(T(\lambda p_1(x)) = \begin{bmatrix}
        \lambda p_1(-1)\\
        \lambda p_1(0) \\
        \lambda p_1(1)
    \end{bmatrix} = \lambda \begin{bmatrix}
        p_1(-1)\\
        p_1(0) \\
        p_1(1)
    \end{bmatrix} = \lambda T(p_1(x)) \)
\end{itemize}

Thus \(T\) is a linear transformation

\subsubsection*{(b)}
Let \(p(x) \in \mathbb{P}_2[x]\), \(p(x) = ax^2 + bx +c\)
\\

\(T(p(x)) = \begin{bmatrix}
    a -b + c \\
    c \\
    a + b +c
\end{bmatrix}\)
\\

Also, \([p(x)]_\mathcal{E} = \begin{bmatrix}
    c \\ 
    b \\ 
    a
\end{bmatrix}\) and \(T(p(x)) = \begin{bmatrix}
    a -b + c \\
    c \\
    a + b +c
\end{bmatrix}\)
\\

So now we need the matrix which maps these together, which we can get via construction
\[
[T]_{\mathcal{E},\mathcal{F}} = \begin{bmatrix}
    1 & -1 & 1 \\
    1 & 0 & 0 \\
    1 & 1 & 1
\end{bmatrix}
\]


\subsection*{Question 6}

For this question we can use the Rank-Nullity Theorem, 
\[
\dim(\ker T) + \dim(\mathrm{Im} \ T) = \dim(V)
\]

Let \(n = \dim(v)\) and \(m = \dim(W)\), we know \(n < m\)

For a linear transformation to be surjective, we need \(\dim(\mathrm{Im} \ T) = m\)

But
\[
\dim(\ker T) + \dim(\mathrm{Im} \ T) = n \\
\]
and since \(\dim(\ker T) \geq 0 \), we have that \(\dim(\mathrm{Im} \ T) \leq n\)
\\

Thus \(\dim(\mathrm{Im} \ T) < m\)

Hence \(T\) cannot be surjective.

\subsection*{Question 7}
Let \(T\) be defined as a linear transformation between the vector spaces \(\mathbb{R}^7\) and 
\(\mathbb{R}^5\)
\[
T: \mathbb{R}^7 \rightarrow \mathbb{R}^5
\]
\(T\) can be represented using a \(5 \times 7\) matrix, \(A\), where 
\(T(\vec{x}) = A\vec{x} \quad \vec{x} \in \mathbb{R}^7\)
\\

For the nullity and rank of \(T\) to be 2 and 5 respectively, we need there to be exactly 5 pivot
columns in \(A\).
\\

One such matrix is 
\[
A = \begin{bmatrix}
    1 & 0 & 0 & 0 & 0 & 0 & 0 \\
    0 & 1 & 0 & 0 & 0 & 0 & 0 \\
    0 & 0 & 1 & 0 & 0 & 0 & 0 \\
    0 & 0 & 0 & 1 & 0 & 0 & 0 \\
    0 & 0 & 0 & 0 & 1 & 0 & 0
\end{bmatrix}
\]


\subsection*{Question 8}
(assuming that the question wants me to show that \(T\) is a linear transformation)
\\

To show that \(T\) is a linear transformation we must show that it preserves addition
and scalar multiplication.
\\

Let \(p_1(x), p_2(x) \in \mathbb{P}[x]\) and let \(\lambda \in \mathbb{R}\)
\begin{itemize}
    \item \(T(p_1(x) + p_2(x)) = \int_{-2}^{1} (p_1(x) + p_2(x)) dx = \int_{-2}^{1} p_1(x) dx
    + \int_{-2}^{1} p_2(x) dx = T(p_1(x)) + T(p_2(x)) \)

    \item \(T(\lambda p_1(x)) = \int_{-2}^{1} (\lambda p_1(x)) dx = 
    \lambda \int_{-2}^{1} p_1(x) dx = \lambda T(p_1(x))\)
\end{itemize}

Thus \(T\) is a linear transformation



\subsection*{Question 9}

\subsubsection*{(a)}
To show that \(T\) is a linear transformation we must show that it preserves addition
and scalar multiplication.
\\

Let \(p_1(x), p_2(x) \in \mathbb{P}_3[x]\) and let \(\lambda \in \mathbb{R}\)
\begin{itemize}
    \item \(T(p_1(x) + p_2(x)) = (3x^2 + x)(p_1 '(x) + p_2 '(x)) = (3x^2 + x)(p_1 '(x))
    + (3x^2 + x)(p_2 '(x)) = T(p_1(x)) + T(p_2(x))\)

    \item \(T(\lambda p_1(x)) = (3x^2 + x)(\lambda p_1 '(x)) = \lambda(3x^2 + x)(p_1 '(x))
    = \lambda T(p_1(x))\)
\end{itemize}

Thus \(T\) is a linear transformation

\subsubsection*{(b)}
\[
\ker T = \{p(x) \in \mathbb{P}_3[x] \mid T(p(x)) = \vec{0}\}
\]

Let \(p(x) \in \mathbb{P}_3[x]\), \(p(x) = ax^3 + bx^2 + cx +d\)

\[T(p(x)) = (3x^2 + x)(3ax^2+2bx +c)\]
\[T(p(x)) = \vec{0} \Rightarrow a = b = c = 0\]

Thus,
\[
\ker T = \{d \in \mathbb{P}_3[x] \mid d \in \mathbb{R}\}
\]

Hence, \(\{1\}\) will for a basis for \(\ker T\)

\subsubsection*{(c)}
\[
\mathrm{Im} \ T = \{ T(p(x)) \mid  p(x) \in \mathbb{P}_3[x]\}
\]

\[
T(p(x)) = (3x^2 + x)(3ax^2+2bx +c) = a(3x^3 + 9x^4) +  b(2x^2 + 6x^3) + c(x + 3x^2)
\]

\[
\Rightarrow \mathrm{Im} \ T = \{a(3x^3 + 9x^4) +  b(2x^2 + 6x^3) + c(x + 3x^2) \mid a,b,c \in \mathbb{R}\}
\]
\[
\mathrm{Im} \ T = \text{span}\{3x^3 + 9x^4, 2x^2 + 6x^3, x + 3x^2\}
\]

This set of vectors forms a basis for the image of \(T\) since we know that the image of \(T\) will be 
of dimension 3 from the rank-nullity theorem.
\[
\dim(\ker T) + \dim(\mathrm{Im} \ T) = \dim(V)
\]
\[
1 + \dim(\mathrm{Im} \ T) = 4
\]
\[
\dim(\mathrm{Im} \ T) = 3
\]

\subsubsection*{(d)}
\[
\dim(\ker T) + \dim(\mathrm{Im} \ T) = \dim(V)
\]
\[
1 + 3 = 4 = 4
\]


\subsection*{Question 10}
We know that the columns of the matrix will be the images of the standard basis vectors
under \(T\), we just need to figure out where these get mapped to.
\\

Figuring this out isn't too hard, using the facts that \(\sin(\theta), \cos(\theta)\) correspond to the 
\(x\) and \(y\) coordinate of a right triangle with unit hyotenuse, which is what you get when you overlay 
the new rotated space over the old space.

\[
(1,0,0) \mapsto (\cos(\theta) ,\sin(\theta) , 0) 
\]
\[
(0,1,0) \mapsto (-\sin(\theta) ,\cos(\theta) , 0)
\]
\[
(0,0,1) \mapsto (0,0,1) 
\]

Thus,
\[
\begin{bmatrix}
    \cos(\theta) & -\sin(\theta) & 0 \\
    \sin(\theta) & \cos(\theta) & 0 \\
    0 & 0 & 1
\end{bmatrix}
\]

\end{document}